\hypertarget{python-3.12-to-3.13-subinterpreters-per-interpreter-gil-parallelism}{%
\chapter{Python 3.12 to 3.13: Subinterpreters \& Per-Interpreter GIL
Parallelism}\label{python-3.12-to-3.13-subinterpreters-per-interpreter-gil-parallelism}}

\hypertarget{isolation-of-interpreter-state-pep-684}{%
\subsection{22.1 Isolation of Interpreter State (PEP
684)}\label{isolation-of-interpreter-state-pep-684}}

Before Python 3.12, multi-threaded concurrency was bottlenecked by the
Global Interpreter Lock (GIL) because all threads shared a single global
interpreter state. PEP 684 introduced support for \textbf{isolated
subinterpreters}, each running with its own per-interpreter GIL. This
configuration allows a single OS process to run multiple Python
interpreters concurrently, achieving true multi-core parallel execution.

\hypertarget{process-memory-layout-shared-vs.-isolated-states}{%
\subsubsection{1. Process Memory Layout: Shared vs.~Isolated
States}\label{process-memory-layout-shared-vs.-isolated-states}}

\begin{lstlisting}
================================ OS PROCESS MEMORY ================================
|                                                                                 |
|  [Global Process State] (Loaded DLLs, static allocation, file descriptors)     |
|                                                                                 |
|  [Subinterpreter 0] (Own GIL)     [Subinterpreter 1] (Own GIL)                 |
|   - Heap: Private Objects          - Heap: Private Objects                     |
|   - Modules: Isolated Dict         - Modules: Isolated Dict                    |
|   - Garbage Collector States       - Garbage Collector States                  |
|   - PyInterpreterState Struct      - PyInterpreterState Struct                 |
|                                                                                 |
===================================================================================
\end{lstlisting}

\hypertarget{the-pyinterpreterstate-structure}{%
\subsubsection{\texorpdfstring{2. The \texttt{PyInterpreterState}
Structure}{2. The PyInterpreterState Structure}}\label{the-pyinterpreterstate-structure}}

To achieve isolation, all mutable structures have been shifted from
global scopes into per-interpreter state wrappers defined in
\passthrough{\lstinline!pycore\_interp.h!}:

\begin{lstlisting}[language=C]
struct _is {
    struct _is *next;                 /* Pointer to next subinterpreter in process */
    struct _ceval_state ceval;        /* Per-interpreter evaluation loop configuration and GIL */
    struct _gc_state gc;              /* Per-interpreter garbage collector generations */
    PyObject *modules;                /* Isolated dictionary containing loaded modules */
    PyObject *sysdict;                /* Isolated sys module context variables */
    /* ... additional fields ... */
};
typedef struct _is PyInterpreterState;
\end{lstlisting}

\begin{center}\rule{0.5\linewidth}{0.5pt}\end{center}

\hypertarget{inter-interpreter-communication-channels}{%
\subsection{22.2 Inter-Interpreter Communication
Channels}\label{inter-interpreter-communication-channels}}

Because subinterpreters share no Python objects directly (to avoid
reference count race conditions across distinct heaps), they cannot
share pointers. Data must be passed between subinterpreters using the
C-level serialization protocols of the
\passthrough{\lstinline!\_xxsubinterpreters!} module.

\hypertarget{channel-mechanics}{%
\subsubsection{1. Channel Mechanics}\label{channel-mechanics}}

\begin{itemize}
\tightlist
\item
  \textbf{Sender}: Serializes/marshals the Python object data into an
  isolated C buffer.
\item
  \textbf{Queue}: Pushes the serialized buffer onto an OS-level
  thread-safe memory queue.
\item
  \textbf{Receiver}: Pops the buffer from the queue, deserializes the
  bytes, and reconstructs a new \passthrough{\lstinline!PyObject!}
  structure on its private heap.
\end{itemize}

This shared-nothing execution model ensures that concurrency remains
lock-free, avoiding heap synchronization bottlenecks across
subinterpreters.

\begin{center}\rule{0.5\linewidth}{0.5pt}\end{center}
